\documentclass[11pt]{article}
% The preamble every project's paper shares: % The preamble every project's paper shares: % The preamble every project's paper shares: \input{../../preamble} from
% projects/<project>/paper/. Boards come from chessboard and xskak, plots from
% pgfplots; every number a paper prints comes from its generated.tex.
\usepackage[LSBC4,T1]{fontenc}
\usepackage[utf8]{inputenc}
\usepackage{lmodern}
\usepackage[margin=1in]{geometry}
\usepackage{microtype}
\usepackage{amsmath}
\usepackage{booktabs}
\usepackage{tabularx}
\usepackage{array}
\usepackage{xcolor}
\usepackage[obeyspaces,spaces]{url}
\usepackage{xskak}
\usepackage{pgfplots}
\pgfplotsset{compat=1.18}
\usepackage{caption}
\usepackage{subcaption}
\usepackage{enumitem}
\usepackage[section]{placeins}
\usepackage[hidelinks]{hyperref}

\setchessboard{boardfontencoding=LSBC4, boardfontfamily=skaknew, blackfieldcolor=black!22,
  whitefieldcolor=white, setfontcolors, showmover=false, label=true, labelfontsize=6pt,
  boardfontsize=14pt}

% A move or line in UCI, and a position in FEN, both breakable at a slash.
\newcommand{\uci}[1]{\texttt{#1}}
\newcommand{\fen}[1]{{\footnotesize\path{#1}}}
\newcolumntype{L}{>{\raggedright\arraybackslash}X}

\setlength{\parskip}{0.4em}
\setlength{\parindent}{0pt}
 from
% projects/<project>/paper/. Boards come from chessboard and xskak, plots from
% pgfplots; every number a paper prints comes from its generated.tex.
\usepackage[LSBC4,T1]{fontenc}
\usepackage[utf8]{inputenc}
\usepackage{lmodern}
\usepackage[margin=1in]{geometry}
\usepackage{microtype}
\usepackage{amsmath}
\usepackage{booktabs}
\usepackage{tabularx}
\usepackage{array}
\usepackage{xcolor}
\usepackage[obeyspaces,spaces]{url}
\usepackage{xskak}
\usepackage{pgfplots}
\pgfplotsset{compat=1.18}
\usepackage{caption}
\usepackage{subcaption}
\usepackage{enumitem}
\usepackage[section]{placeins}
\usepackage[hidelinks]{hyperref}

\setchessboard{boardfontencoding=LSBC4, boardfontfamily=skaknew, blackfieldcolor=black!22,
  whitefieldcolor=white, setfontcolors, showmover=false, label=true, labelfontsize=6pt,
  boardfontsize=14pt}

% A move or line in UCI, and a position in FEN, both breakable at a slash.
\newcommand{\uci}[1]{\texttt{#1}}
\newcommand{\fen}[1]{{\footnotesize\path{#1}}}
\newcolumntype{L}{>{\raggedright\arraybackslash}X}

\setlength{\parskip}{0.4em}
\setlength{\parindent}{0pt}
 from
% projects/<project>/paper/. Boards come from chessboard and xskak, plots from
% pgfplots; every number a paper prints comes from its generated.tex.
\usepackage[LSBC4,T1]{fontenc}
\usepackage[utf8]{inputenc}
\usepackage{lmodern}
\usepackage[margin=1in]{geometry}
\usepackage{microtype}
\usepackage{amsmath}
\usepackage{booktabs}
\usepackage{tabularx}
\usepackage{array}
\usepackage{xcolor}
\usepackage[obeyspaces,spaces]{url}
\usepackage{xskak}
\usepackage{pgfplots}
\pgfplotsset{compat=1.18}
\usepackage{caption}
\usepackage{subcaption}
\usepackage{enumitem}
\usepackage[section]{placeins}
\usepackage[hidelinks]{hyperref}

\setchessboard{boardfontencoding=LSBC4, boardfontfamily=skaknew, blackfieldcolor=black!22,
  whitefieldcolor=white, setfontcolors, showmover=false, label=true, labelfontsize=6pt,
  boardfontsize=14pt}

% A move or line in UCI, and a position in FEN, both breakable at a slash.
\newcommand{\uci}[1]{\texttt{#1}}
\newcommand{\fen}[1]{{\footnotesize\path{#1}}}
\newcolumntype{L}{>{\raggedright\arraybackslash}X}

\setlength{\parskip}{0.4em}
\setlength{\parindent}{0pt}

% Written by projects/maia-rating-offset/paper/generate.mjs; do not edit by hand.

\newcommand{\AppOffset}{250}
\newcommand{\OffsetMin}{0}
\newcommand{\OffsetMax}{500}
\newcommand{\OffsetStep}{25}
\newcommand{\OffsetCount}{21}
\newcommand{\MaiaStart}{1100}
\newcommand{\MaiaEnd}{2000}
\newcommand{\MaiaWidth}{100}
\newcommand{\MaiaBands}{11}
\newcommand{\AppCommit}{\texttt{3eb8afd}}
\newcommand{\LowestBelow}{850}
\newcommand{\AppOpponentOffset}{400}
\newcommand{\MinPly}{11}
\newcommand{\MinClock}{30}
\newcommand{\Games}{656}
\newcommand{\Moves}{34{,}025}
\newcommand{\HimMoves}{17{,}016}
\newcommand{\OppMoves}{17{,}009}
\newcommand{\FitGames}{328}
\newcommand{\FitMoves}{17{,}242}
\newcommand{\HeldGames}{328}
\newcommand{\HeldMoves}{16{,}783}
\newcommand{\RatingMin}{513}
\newcommand{\RatingMedian}{1017}
\newcommand{\RatingMax}{1228}
\newcommand{\MedianBand}{1200}
\newcommand{\Draws}{1{,}000}
\newcommand{\Seed}{7}
\newcommand{\TenMinuteGames}{655}
\newcommand{\OtherGames}{one}
\newcommand{\Fitted}{250}
\newcommand{\Kept}{250}
\newcommand{\HeldDiff}{0.0000}
\newcommand{\HeldLo}{0.0000}
\newcommand{\HeldHi}{0.0000}
\newcommand{\FitGain}{0.0000}
\newcommand{\LLBase}{$-$1.393}
\newcommand{\LLFitted}{$-$1.393}
\newcommand{\TopBase}{54.0}
\newcommand{\TopFitted}{54.0}
\newcommand{\BestAll}{250}
\newcommand{\BestAllLo}{250}
\newcommand{\BestAllHi}{275}
\newcommand{\BestHim}{175}
\newcommand{\BestHimLo}{150}
\newcommand{\BestHimHi}{200}
\newcommand{\BestOpp}{425}
\newcommand{\BestOppLo}{350}
\newcommand{\BestOppHi}{425}
\newcommand{\HeldBest}{250}
\newcommand{\LossAtMin}{0.0213}
\newcommand{\LossAtMax}{0.0142}
\newcommand{\RatioAtMin}{2.2}
\newcommand{\RatioAtMax}{1.4}
\newcommand{\TopMin}{53.5}
\newcommand{\TopMax}{54.0}
\newcommand{\GainHim}{0.0018}
\newcommand{\GainOpp}{0.0064}
\newcommand{\GainLow}{0.0038}
\newcommand{\GainHigh}{0.0001}
\newcommand{\CurveFit}{(0,-0.02236) (25,-0.01693) (50,-0.01423) (75,-0.01259) (100,-0.00995) (125,-0.00770) (150,-0.00563) (175,-0.00367) (200,-0.00249) (225,-0.00169) (250,0.00000) (275,-0.00006) (300,-0.00069) (325,-0.00061) (350,-0.00061) (375,-0.00198) (400,-0.00305) (425,-0.00466) (450,-0.00634) (475,-0.00909) (500,-0.01244)}
\newcommand{\CurveHeld}{(0,-0.02030) (25,-0.01678) (50,-0.01391) (75,-0.01129) (100,-0.00911) (125,-0.00679) (150,-0.00466) (175,-0.00208) (200,-0.00132) (225,-0.00103) (250,0.00000) (275,-0.00013) (300,-0.00109) (325,-0.00131) (350,-0.00199) (375,-0.00289) (400,-0.00508) (425,-0.00617) (450,-0.00868) (475,-0.01153) (500,-0.01607)}
\newcommand{\CurveAll}{(0,-0.02134) (25,-0.01685) (50,-0.01407) (75,-0.01195) (100,-0.00954) (125,-0.00725) (150,-0.00515) (175,-0.00289) (200,-0.00191) (225,-0.00136) (250,0.00000) (275,-0.00009) (300,-0.00089) (325,-0.00096) (350,-0.00129) (375,-0.00243) (400,-0.00405) (425,-0.00540) (450,-0.00749) (475,-0.01030) (500,-0.01423)}
\newcommand{\CurveHim}{(0,-0.00792) (25,-0.00461) (50,-0.00291) (75,-0.00184) (100,-0.00095) (125,0.00018) (150,0.00099) (175,0.00177) (200,0.00125) (225,0.00022) (250,0.00000) (275,-0.00186) (300,-0.00436) (325,-0.00643) (350,-0.00789) (375,-0.01078) (400,-0.01408) (425,-0.01719) (450,-0.02064) (475,-0.02522) (500,-0.03109)}
\newcommand{\CurveOpp}{(0,-0.03478) (25,-0.02910) (50,-0.02524) (75,-0.02207) (100,-0.01812) (125,-0.01468) (150,-0.01130) (175,-0.00755) (200,-0.00507) (225,-0.00295) (250,0.00000) (275,0.00167) (300,0.00259) (325,0.00452) (350,0.00531) (375,0.00593) (400,0.00597) (425,0.00639) (450,0.00566) (475,0.00463) (500,0.00263)}
\newcommand{\CurveLow}{(0,-0.00833) (25,-0.00833) (50,-0.00833) (75,-0.00833) (100,-0.00842) (125,-0.00842) (150,-0.00800) (175,-0.00525) (200,-0.00346) (225,-0.00289) (250,0.00000) (275,0.00058) (300,0.00101) (325,0.00229) (350,0.00314) (375,0.00354) (400,0.00384) (425,0.00381) (450,0.00170) (475,-0.00066) (500,-0.00257)}
\newcommand{\LowFrom}{513}
\newcommand{\LowTo}{936}
\newcommand{\LowMoves}{11{,}296}
\newcommand{\LowBest}{400}
\newcommand{\LowLo}{325}
\newcommand{\LowHi}{425}
\newcommand{\CurveMid}{(0,-0.03016) (25,-0.02770) (50,-0.02095) (75,-0.01675) (100,-0.01284) (125,-0.00913) (150,-0.00398) (175,-0.00124) (200,-0.00081) (225,-0.00082) (250,0.00000) (275,-0.00092) (300,-0.00156) (325,-0.00235) (350,-0.00406) (375,-0.00698) (400,-0.00838) (425,-0.01026) (450,-0.01388) (475,-0.01762) (500,-0.01947)}
\newcommand{\MidFrom}{937}
\newcommand{\MidTo}{1077}
\newcommand{\MidMoves}{11{,}263}
\newcommand{\MidBest}{250}
\newcommand{\MidLo}{200}
\newcommand{\MidHi}{250}
\newcommand{\CurveHigh}{(0,-0.02551) (25,-0.01459) (50,-0.01298) (75,-0.01080) (100,-0.00739) (125,-0.00426) (150,-0.00350) (175,-0.00219) (200,-0.00147) (225,-0.00039) (250,0.00000) (275,0.00005) (300,-0.00209) (325,-0.00279) (350,-0.00294) (375,-0.00384) (400,-0.00758) (425,-0.00970) (450,-0.01028) (475,-0.01260) (500,-0.02058)}
\newcommand{\HighFrom}{1078}
\newcommand{\HighTo}{1228}
\newcommand{\HighMoves}{11{,}466}
\newcommand{\HighBest}{275}
\newcommand{\HighLo}{225}
\newcommand{\HighHi}{275}
\newcommand{\OffsetTable}{%
\begin{tabular}{lrrrr}
\toprule
offset & fitting half & held-out half & all games & first choice \\
\midrule
0 & $-$1.4165 & $-$1.4120 & $-$1.4143 & 53.5\% \\
25 & $-$1.4110 & $-$1.4085 & $-$1.4098 & 53.6\% \\
50 & $-$1.4083 & $-$1.4057 & $-$1.4070 & 53.7\% \\
75 & $-$1.4067 & $-$1.4030 & $-$1.4049 & 53.7\% \\
100 & $-$1.4040 & $-$1.4008 & $-$1.4025 & 53.7\% \\
125 & $-$1.4018 & $-$1.3985 & $-$1.4002 & 53.8\% \\
150 & $-$1.3997 & $-$1.3964 & $-$1.3981 & 53.8\% \\
175 & $-$1.3978 & $-$1.3938 & $-$1.3958 & 53.9\% \\
200 & $-$1.3966 & $-$1.3931 & $-$1.3948 & 53.9\% \\
225 & $-$1.3958 & $-$1.3928 & $-$1.3943 & 53.9\% \\
250 (app, fitted) & $-$1.3941 & $-$1.3917 & $-$1.3929 & 54.0\% \\
275 & $-$1.3942 & $-$1.3919 & $-$1.3930 & 54.0\% \\
300 & $-$1.3948 & $-$1.3928 & $-$1.3938 & 54.0\% \\
325 & $-$1.3947 & $-$1.3931 & $-$1.3939 & 54.0\% \\
350 & $-$1.3947 & $-$1.3937 & $-$1.3942 & 54.0\% \\
375 & $-$1.3961 & $-$1.3946 & $-$1.3954 & 53.9\% \\
400 & $-$1.3971 & $-$1.3968 & $-$1.3970 & 53.9\% \\
425 & $-$1.3987 & $-$1.3979 & $-$1.3983 & 53.8\% \\
450 & $-$1.4004 & $-$1.4004 & $-$1.4004 & 53.8\% \\
475 & $-$1.4032 & $-$1.4033 & $-$1.4032 & 53.7\% \\
500 & $-$1.4065 & $-$1.4078 & $-$1.4072 & 53.7\% \\
\bottomrule
\end{tabular}
}
\newcommand{\LookTable}{%
\begin{tabular}{lrrl}
\toprule
moves & count & best offset & 95\% interval over draws \\
\midrule
all & 34{,}025 & 250 & 250 to 275 \\
the player's & 17{,}016 & 175 & 150 to 200 \\
his opponents' & 17{,}009 & 425 & 350 to 425 \\
movers rated 513 to 936 & 11{,}296 & 400 & 325 to 425 \\
movers rated 937 to 1077 & 11{,}263 & 250 & 200 to 250 \\
movers rated 1078 to 1228 & 11{,}466 & 275 & 225 to 275 \\
\bottomrule
\end{tabular}
}
\newcommand{\ExampleRating}{1061}
\newcommand{\ExampleOppRating}{1082}
\newcommand{\ExamplePlayed}{29.~\symking{}h1}
\newcommand{\ExampleLowP}{41}
\newcommand{\ExampleHighP}{95}
\newcommand{\Example}{% Written by projects/maia-rating-offset/paper/generate.mjs; do not edit by hand.
\noindent\begin{minipage}[t]{0.37\textwidth}\vspace{0pt}
\chessboard[setfen={4r1rk/1R4pp/p1Q5/8/5P2/4q3/P5PP/5RK1 w - - 0 29}, boardfontsize=12pt, pgfstyle=straightmove, arrow=to, linewidth=0.07em, color=red!75!black, markmoves={g1-h1}]
\end{minipage}\hfill
\begin{minipage}[t]{0.6\textwidth}\vspace{0pt}\small\setlength{\tabcolsep}{4pt}
\textbf{Played:} 29.~\symking{}h1 (red), by the player at 1061 against 1082.\par\smallskip
\begin{tabular}{rllrl}
\toprule
offset & his band & opponent's & played & Maia-2's first choice \\
\midrule
0 & <1100 & <1100 & 40.6\% & 29.~\symrook{}f2, 59.4\% \\
250 & 1300 & 1300 & 83.1\% & 29.~\symking{}h1, 83.1\% \\
500 & 1500 & 1500 & 95.2\% & 29.~\symking{}h1, 95.2\% \\
\bottomrule
\end{tabular}\par\smallskip
\textbf{UCI:} played \uci{g1h1}; first choices \uci{f1f2}, \uci{g1h1}, \uci{g1h1}\par
\textbf{FEN:} \fen{4r1rk/1R4pp/p1Q5/8/5P2/4q3/P5PP/5RK1 w - - 0 29}
\end{minipage}%
}


\title{Which rating does Maia-2 need for this player?\\[0.3em]
\large Fitting the chess.com to Lichess offset on the player's own games}
\author{Chess Coach research lab}
\date{}

\begin{document}
\maketitle

\begin{abstract}
The coaching app converts the player's chess.com rating to Maia-2's Lichess scale by adding
\AppOffset{} points, a number its own code calls a default, not a measurement. We measured it on
the \Games{} rapid games of the player's archive: at every offset from \OffsetMin{} to
\OffsetMax, how likely Maia-2 finds the \Moves{} moves both players made. By a rule fixed before
any number was computed, the fitting half of the games chose \Fitted, the app's own number, so it
stays; the held-out half chose it too. Beside the verdict, a single offset hides a split: Maia-2
predicts the player's own moves best at \BestHim{} and his opponents' at \BestOpp, with intervals
that do not overlap. Whether separate offsets beat one is a question for its own rule.
\end{abstract}

\section{The question}

Maia-2 predicts the move a human plays from the position and two ratings, the mover's and the
opponent's. It learned from Lichess games, so its ratings are Lichess ratings, in \MaiaBands{}
bands: one below \MaiaStart, one for every \MaiaWidth{} points up to \MaiaEnd, and one from
\MaiaEnd{} up. The coaching app knows the player's chess.com rating and adds \AppOffset{} points
to it before it asks Maia-2 anything; its own code calls that number a default, not a
measurement. Every Maia number the app gives rests on it: the odds that a mistake gets punished,
how findable a move is at the player's level, and the Maia bot's model of the player and of the
opponent it plays as.

The player's games can measure it. The offset Maia-2 needs is the one under which it finds the
moves actually played most likely, with both players at their own chess.com ratings converted by
that offset.

\section{The moves}

The games are the \Games{} in the app's archive, every one from the standard start with both
ratings and a clock after every ply. All are rapid: \TenMinuteGames{} at ten minutes a side and
\OtherGames{} at a slower control. A move counts from ply \MinPly{} on when both players had at least
\MinClock{} seconds left at the position, the protocol of the Maia papers that the lab's Maia-2
benchmark also follows. That leaves \Moves{} moves: \HimMoves{} by the player and \OppMoves{} by
his opponents. The movers are rated from \RatingMin{} to \RatingMax{} on chess.com (median
\RatingMedian, which the app's offset puts in Maia's \MedianBand{} band); the rating is the one
chess.com recorded after each game.

\section{The instrument and the rule}

At an offset $o$, a move by a player rated $r$ against one rated $r'$ scores the logarithm of
Maia-2's probability of the move played, with the mover's band the one $r + o$ falls in and the
opponent's the one $r' + o$ falls in, exactly as the app computes them (\texttt{defaultMaiaElo}),
through the app's own Maia-2 in Node (\texttt{humanPolicy} with both bands). The offsets run over
the range the app allows, \OffsetMin{} to \OffsetMax, in steps of \OffsetStep{} (\OffsetCount{}
offsets). The bands are \MaiaWidth{} points wide, so a step changes only the players whose
converted rating crosses a band edge, and at the app's offset everyone rated under
\LowestBelow{} on chess.com shares the lowest band.

The rule was committed before any number was computed. The games, ordered by when they ended,
alternate between a fitting half (\FitGames{} games, \FitMoves{} moves) and a held-out half
(\HeldGames{} games, \HeldMoves{} moves). The fitted offset is the one with the highest mean
log-likelihood per move on the fitting half. It replaces \AppOffset{} only if, on the held-out
half, its mean log-likelihood per move beats \AppOffset's with the 95\% interval of the
difference, over \Draws{} bootstrap draws of the held-out games, entirely above zero.

\section{Results}

The fitting half's best offset is \Fitted, the app's own, so \AppOffset{} stays. The held-out
half's best is \HeldBest{} as well, and so is the best over all the games; over \Draws{}
bootstrap draws of the games, the best offset falls between \BestAllLo{} and \BestAllHi{}
(Table~\ref{tab:offsets}, Figure~\ref{fig:curve}).

Near the top the offset matters little. At \AppOffset{} the moves played average a
log-likelihood of \LLBase{} per move; \OffsetMin{} loses \LossAtMin{} per move and \OffsetMax{}
loses \LossAtMax. As a geometric mean, Maia-2 gives the moves played \RatioAtMin\% more
probability at \AppOffset{} than at \OffsetMin, and \RatioAtMax\% more than at \OffsetMax. Its
first choice is the move played \TopMin\% to \TopMax\% of the time across the whole range: the
offset changes how sure Maia-2 is far more than which move it ranks first.

\begin{figure}[!htbp]
\centering
\begin{tikzpicture}
\begin{axis}[width=0.7\textwidth, height=0.42\textwidth, xmin=-10, xmax=510,
  xlabel={offset added to the chess.com rating}, ylabel={change from \AppOffset, per move},
  xtick={0,100,200,300,400,500}, grid=major, legend pos=south east, legend cell align=left,
  y tick label style={/pgf/number format/fixed, /pgf/number format/precision=3},
  scaled y ticks=false]
\addplot[black, mark=*, mark size=1.2pt] coordinates {\CurveFit};
\addlegendentry{fitting half}
\addplot[red!70!black, dashed, mark=square*, mark size=1.2pt] coordinates {\CurveHeld};
\addlegendentry{held-out half}
\end{axis}
\end{tikzpicture}
\caption{The mean log-likelihood per move of the moves played at each offset, as a change from
its value at the app's \AppOffset. Both halves peak at \AppOffset.}
\label{fig:curve}
\end{figure}

\begin{table}[!htbp]
\centering\small
\OffsetTable
\caption{The mean log-likelihood per move at each offset, on each half and over all the games,
and the share of all the moves that are Maia-2's first choice.}
\label{tab:offsets}
\end{table}

\section{Beside the verdict}

Named before the first run, these looks decide nothing, but they show what one offset hides
(Table~\ref{tab:looks}, Figure~\ref{fig:looks}). Maia-2 predicts the player's own moves best at
\BestHim{} (\BestHimLo{} to \BestHimHi{} over the draws) and his opponents' best at \BestOpp{}
(\BestOppLo{} to \BestOppHi); the two intervals do not overlap. By the movers' rating, the
weakest third (\LowFrom{} to \LowTo) fits best at \LowBest{} (\LowLo{} to \LowHi), an interval
that leaves \AppOffset{} out; the middle third fits best at \MidBest{} (\MidLo{} to \MidHi) and
the strongest at \HighBest{} (\HighLo{} to \HighHi), and both intervals hold \AppOffset, like
the moves as a whole. The gains a part's own best offset would bring over \AppOffset{} are
small: \GainHim{} per move for the player's moves, \GainOpp{} for his opponents', \GainLow{}
for the weakest third and \GainHigh{} for the strongest.

One offset moves both players at once, so it cannot follow the split: every offset that fits his
opponents better than \AppOffset{} fits him worse, and every one that fits him better fits them
worse. \AppOffset{} is where the two balance. Whether an offset for the player and another for
his opponents would beat one on held-out games is a question for its own project, with its rule
written first: these looks chose where to look, so they cannot also be the test. A split would
also move both directions of the punish odds, which the punish-odds project found calibrated at
\AppOffset, so it would have to keep them so.

\begin{table}[!htbp]
\centering\small
\LookTable
\caption{The best offset for each part of the moves, and its 95\% interval over \Draws{}
bootstrap draws of the games. The thirds split the movers' chess.com ratings.}
\label{tab:looks}
\end{table}

\begin{figure}[!htbp]
\centering
\begin{tikzpicture}
\begin{axis}[name=sides, width=0.48\textwidth, height=0.4\textwidth, xmin=-10, xmax=510,
  xlabel={offset}, ylabel={change from \AppOffset, per move},
  xtick={0,250,500}, grid=major, legend pos=south west, legend cell align=left,
  legend style={font=\scriptsize}, y tick label style={/pgf/number format/fixed,
  /pgf/number format/precision=3}, scaled y ticks=false, title={by side}]
\addplot[black, mark=*, mark size=1pt] coordinates {\CurveHim};
\addlegendentry{the player's moves}
\addplot[red!70!black, dashed, mark=square*, mark size=1pt] coordinates {\CurveOpp};
\addlegendentry{his opponents'}
\end{axis}
\begin{axis}[at={(sides.east)}, anchor=west, xshift=1cm, width=0.48\textwidth,
  height=0.4\textwidth, xmin=-10, xmax=510, xlabel={offset}, xtick={0,250,500}, grid=major,
  legend pos=south west, legend cell align=left, legend style={font=\scriptsize},
  y tick label style={/pgf/number format/fixed, /pgf/number format/precision=3},
  scaled y ticks=false, title={by the mover's rating}]
\addplot[blue!70!black, mark=triangle*, mark size=1.2pt] coordinates {\CurveLow};
\addlegendentry{\LowFrom--\LowTo}
\addplot[black, mark=*, mark size=1pt] coordinates {\CurveMid};
\addlegendentry{\MidFrom--\MidTo}
\addplot[gray, dashed, mark=o, mark size=1pt] coordinates {\CurveHigh};
\addlegendentry{\HighFrom--\HighTo}
\end{axis}
\end{tikzpicture}
\caption{The same curve for each part of the moves, as a change from each part's value at
\AppOffset.}
\label{fig:looks}
\end{figure}

Figure~\ref{fig:example} shows what the bands do in one position: the player's move whose
probability moves most between \OffsetMin{} and \OffsetMax. The stronger the bands the offset
gives the two players, the surer Maia-2 is of the move he chose.

\begin{figure}[!htbp]\Example\caption{The player's move at the bands three offsets give the two
players: Maia-2's probability of the move played and its first choice.}\label{fig:example}
\end{figure}

\section{What it means for the app}

Nothing to port. The app's \AppOffset{} is the offset this player's games choose, so the punish
odds, how findable a move is and the Maia bot stand as they are. The app lets the player change
the offset; any offset from \BestAllLo{} to \BestAllHi, the range the best one takes over the
draws, fits these games about as well. The split between the player and his opponents is the next question for this
part of the app.

\section{Limitations}

\begin{itemize}[nosep]
\item One player's games and his opponents', rated \RatingMin{} to \RatingMax{} on chess.com. The
offset that fits them says nothing about players far from that range.
\item The rating is chess.com's after each game, a few points from the one the game was played
at.
\item Maia-2's lowest band holds everyone under \MaiaStart{} on its scale, so it cannot tell the
weakest players here apart from each other at any offset that leaves them there.
\item The likelihood measures how well Maia-2 predicts the moves, not why. An offset that fits
best also absorbs whatever Maia-2 gets wrong about this pool of players, which is what the app
needs from it, but it is not a statement about the two sites' rating pools.
\end{itemize}

\end{document}
